% Generated draft: requires mathematical and source review.
\documentclass[11pt]{article}
\usepackage[margin=25mm]{geometry}
\usepackage{fontspec}
\setmainfont{Latin Modern Roman}
\usepackage{amsmath,amssymb,mathtools,unicode-math}
\setmathfont{Latin Modern Math}
\usepackage{enumitem,needspace,xurl,hyperref}
\setlength{\parindent}{0pt}
\setlength{\parskip}{5pt}
\setlength{\emergencystretch}{3em}
\newcommand{\N}{\mathbb{N}}
\newcommand{\Z}{\mathbb{Z}}
\newcommand{\Q}{\mathbb{Q}}
\newcommand{\R}{\mathbb{R}}
\newcommand{\C}{\mathbb{C}}
\newcommand{\F}{\mathbb{F}}
\newcommand{\A}{\mathbb{A}}
\newcommand{\PP}{\mathbb{P}}
\newcommand{\E}{\mathbb{E}}
\newcommand{\eps}{\varepsilon}
\newcommand{\dd}{\,\mathrm{d}}
\begin{document}
\section*{UnsolvedMath ID: 800822}
\textbf{Source problem code:} catalogue entry
\subsection{Definitions and mathematical statement}
A conventional setting for the two-dimensional incompressible Euler equation is a smooth oriented Riemannian surface \((M,g)\), with Riemannian area form \(\mu\), a time-dependent velocity field \(u(t,\cdot)\), and a pressure \(p(t,\cdot)\), satisfying
\[
\partial_t u+\nabla_u u=-\operatorname{grad}_g p,
\qquad \operatorname{div}_{\mu}u=0.
\]
Here \(\nabla\) is the Levi--Civita connection of \(g\). This supplies conventional notation only: the record does not specify the surface, metric, regularity class, boundary conditions, or time interval intended in the problem.

To describe the supplied enstrophy background in a standard setting, suppose additionally that \(M\) is compact without boundary and that the fields are smooth. Write \(\alpha=u^{\flat}\), where \(u^{\flat}=g(u,\cdot)\), and define the scalar vorticity \(\omega\) by
\[
\mathrm d\alpha=\omega\mu.
\]
The kinetic energy and generalized enstrophy functionals are
\[
H_g(u)=\frac12\int_M g(u,u)\,\mu,
\qquad
C_F([\alpha])=\int_M F(\omega)\,\mu
\quad\bigl(F\in C^{\infty}(\R,\R)\bigr).
\]
Here \([\alpha]\) denotes the class of \(\alpha\) modulo exact one-forms. The enstrophy functionals are Casimirs in the usual Lie--Poisson description: they are constant on coadjoint orbits of the group of area-preserving diffeomorphisms. With the area form fixed, their expression in terms of \([\alpha]\) does not require a Riemannian metric; the identification of velocity with \(\alpha\), and the kinetic energy, do involve the metric. These additional compactness, orientation, and boundary assumptions make this explanatory framework explicit; they are not hypotheses supplied for the requested problem.

\textbf{Requested statement, as supplied.} Prove the non-integrability of the two-dimensional Euler equation.

\textbf{Unresolved qualifications.} The record does not define integrability or specify the class of admissible first integrals, the phase space, or whether the assertion concerns the full equation, a particular coadjoint orbit, or a generic class of orbits. Consequently, it does not determine a unique precise non-integrability conjecture. The source describes the enstrophy Casimirs as specifying a coadjoint orbit, namely a set of isovortical vector fields, while providing no classification hypotheses or additional orbit data. This description is retained as source background, without asserting that generalized enstrophy values alone completely classify coadjoint orbits. The supplied distinction is that these Casimirs constrain the orbit on which evolution takes place but do not describe the dynamics within that orbit.
\subsection{Short English statement}
The supplied problem asks for a proof that the two-dimensional incompressible Euler equation is non-integrable. Its background emphasizes that infinitely many enstrophy Casimirs constrain coadjoint orbits without describing motion within them. The intended meaning of integrability, the surface, and the analytic and boundary settings are unspecified, so a unique precise conjecture requires further formulation. The registry labels the problem open; that label is supplied metadata, not an independently verified assessment.
\subsection{Sources}
\begin{itemize}
\item Local UnsolvedMath registry, permanent ID \(800822\), exact source problem code \texttt{catalogue entry}; supplied title: \emph{Geometric Hydrodynamics in Open Problems --- Problem 21}. The record attributes the problem to Khesin, Misiolek, and Modin, with proposed year 2022, and labels its status \texttt{open}.
\item Khesin, Misiolek, Modin, \emph{Geometric Hydrodynamics in Open Problems} (2022), Problem 21. Supplied source URL: \url{https://arxiv.org/abs/2205.01143}. The supplied provenance records access on 2026-10-16 and extraction by \texttt{arxiv\_latex}. These are supplied provenance details; the source was not independently visited or verified for this draft.
\item Supplied upstream revision: \texttt{d371ee917ac9647fbb06cd6fc529deec6f3dda82}. The local record was created and updated on 2026-10-16. These registry dates do not constitute an independent literature assessment.
\end{itemize}
\end{document}
